% ===================================================================
%  TABULKA č. 1 — Škola. — Lyceum. — Třída.
%
%  Ceci n'est PAS une transcription : c'est une traduction, faite en
%  2026, en tcheque standard. D'ou trois differences avec texto/io et
%  texto/fr, et trois seulement :
%
%   * pas de \nl ni de \cc ;
%   * pas de \begin{VUpage} ;
%   * le gras se pose sur le TERME seul, ponctuation dehors.
%
%  Les cles « %%K » sont celles de texto/io, macro pour macro pour
%  l'apparat, et les renvois \textsuperscript{(n)} visent les MEMES
%  objets de la planche.
%
%  TREIZIEME COLONNE DES DIX-SEPT, deuxieme langue slave apres
%  l'ukrainien, et la premiere du volume qui decline SEPT cas.
%
%  CELA NE COUTERA RIEN AUX RENVOIS, ET C'EST LE POINT A TENIR DES
%  MAINTENANT. Le tcheque n'a pas d'article, mais il POSTPOSE son
%  genitif comme l'ido — « dveře hradu », la porte du chateau — de
%  sorte que la source ordinaire des inversions, celle qui a coute
%  vingt-neuf tours au basque et treize au finnois, ne joue pas ici.
%  L'ukrainien avait donne le meme resultat pour la meme raison.
%
%  CE QUI COUTE, EN REVANCHE, C'EST LE GRAS. Le mot decline change de
%  finale selon sa fonction, et le releve des objets doit porter la
%  forme QUE LA PHRASE DEMANDE, non la forme du dictionnaire : « na
%  tabuli » et non « tabule ». Meme parti qu'en russe, en ukrainien
%  et en basque ; c'est la phrase qui commande, non le lexique.
%
%  ET LE DANGER DE CETTE COLONNE N'EST NI LE RUSSE NI LE POLONAIS :
%  c'est le SLOVAQUE. Les deux langues s'ecrivent presque de la meme
%  facon et se lisent l'une l'autre sans peine ; ce qui les separe a
%  l'oeil tient a une poignee de signes — ä, ô, ľ, ĺ, ŕ — que le
%  tcheque n'ecrit jamais. Le balayage les cherche sur le code du
%  caractere, comme il cherchait les cedilles turques en roumain.
%
%  ET LES NOMS DES NOTES RESTENT CEUX DE L'IDO. Le tcheque appelle H
%  ce que l'ido nomme B ; mais la ligne du fac-simile n'enseigne pas
%  le solfege tcheque, elle donne l'equivalence que l'ido pose, et
%  les onze colonnes precedentes l'ont toutes gardee telle quelle.
%  Une traduction rend ce que le texte DIT, non ce que le lecteur
%  aurait dit a sa place.
%
%  LA CONJONCTION LIVRESQUE ET LE PRESENTATIF BIBLIQUE, otes en 2026.
%  « nebot » revenait trente et une fois : ce n'est pas une faute, le
%  tcheque ecrit l'emploie encore, mais c'est la conjonction des
%  livres, et le livret raconte une planche murale. Elle devient
%  « protoze », qui est la conjonction de tout le monde, et « vzdyt »
%  aux quatre endroits ou la cause est une evidence qu'on rappelle,
%  non une explication qu'on donne.
%
%  « hle » ne revenait que neuf fois — la ou l'espagnol en avait
%  vingt-deux et le galicien trente-trois. Le releveur tcheque avait
%  deja resiste au « voici » du fac-simile ; il n'avait cede que
%  devant « voici que », sept fois sur neuf. C'est bien le « voici
%  que » qui est le plus contagieux, comme en galicien et en
%  neerlandais : « a hle, jsme u », « a hle, vraceji se ». On y met
%  l'adverbe de temps que le francais y mettait : « a uz », « a ted ».
%
%  ET UNE REPETITION QUE LA CORRECTION A MISE AU JOUR. « Je velmi
%  zima, nebot vladne zima » devenait « Je velmi zima, protoze je
%  zima » — le meme mot deux fois dans cinq. Refait en « Je velka
%  zima, protoze prisel mraz ». La tournure vieillie cachait la
%  redite ; l'oter la montre.
% ===================================================================

%%K t01-apar-0 apar
\VUsaut{2.0mm}
\VUcentre{12.6pt}{120}{VYSVĚTLUJÍCÍ KNÍŽKA}
\VUsaut{3.2mm}
\VUcentre{8.4pt}{160}{K}
\VUsaut{2.6mm}
\VUtitre{15.6pt}{0}{Delmasovým pomocným tabulkám}
\VUsaut{3.4mm}
\VUornamento{9pt}
\VUsaut{5.0mm}
\VUcentre{13.2pt}{90}{PRVNÍ SÉRIE}
\VUsaut{3.0mm}
\VUfilet{16mm}
\VUsaut{5.6mm}

%%K t01-apar-1 apar
\VUtitre{15.0pt}{60}{TABULKA č. 1}
\VUsaut{1.4mm}
\VUtitre{11.4pt}{0}{Škola. --- Lyceum. --- Třída.}
\VUsaut{2.2mm}
\VUfilet{20mm}
\VUsaut{6.4mm}

%%K t01-tit-1 sub
\VUpk{10.2pt}{40}{Obecný popis.}
\VUsaut{3.0mm}

%%K t01-01-1 p
1. --- Tato tabulka znázorňuje \VUgras{třídu} v \VUgras{lyceu}. V této třídě je osm
\VUgras{stolů} \textsuperscript{(18)} a osm \VUgras{lavic} \textsuperscript{(32)}. Na
každé lavici sedí tři žáci. \VUgras{Učitel} \textsuperscript{(7)} sedí na
\VUgras{židli} \textsuperscript{(8)} u své \VUgras{katedry} \textsuperscript{(6)}.

%%K t01-02-1 p
2. --- Vlevo od katedry je \VUgras{prosklená skříň} \textsuperscript{(15)}. Vpravo je
\VUgras{černá tabule} \textsuperscript{(1)} se svým \VUgras{hadrem} \textsuperscript{(2)} a
\VUgras{křídou} \textsuperscript{(3)}.

%%K t01-02-2 p
Nad katedrou visí \VUgras{nástěnné tabule} \textsuperscript{(9, 11, 12)} a
\VUgras{mapa} \textsuperscript{(10)}.

%%K t01-03-1 p
3. --- Ne všichni žáci sedí na svém \VUgras{místě} \textsuperscript{(17)}. Jan
\textsuperscript{(4)} stojí u tabule; drží \VUgras{ten} hadr a maže
\VUgras{geometrické obrazce.}

%%K t01-03-2 p
Petr je u katedry; sbírá \VUgras{písemné úkoly} \textsuperscript{(5)} a nese je učiteli.

%%K t01-04-1 p
4. --- \VUgras{Tento} učitel je učitelem \VUgras{živých jazyků.} Učí žáky mluvit, číst a psát
cizími jazyky \VUgras{(německy, anglicky, španělsky, italsky, rusky} nebo
\VUgras{francouzsky)}.

%%K t01-05-1 p
5. --- Třída je velmi prostorná a velmi světlá. Vzadu je široké \VUgras{okno} se dvěma
\VUgras{okenicemi} \textsuperscript{(78)} a \VUgras{prosklené dveře,} aby se dala větrat a
osvětlovat.

%%K t01-05-2 p
V této třídě není zima. Uprostřed stojí, aby ji vytápěla, velmi dobrá \VUgras{kamna}
\textsuperscript{(46)} se svou \VUgras{rourou} \textsuperscript{(45)}.

%%K t01-05-3 p
Na stěně jsou připevněné \VUgras{věšáky} \textsuperscript{(58)}; visí na nich čepice a pláště
školáků.

%%K t01-tit-2 sub
\VUsaut{5.2mm}
\VUpk{10.2pt}{40}{Školní dvůr.}
\VUsaut{3.6mm}

%%K t01-06-1 p
6. --- \VUgras{Hodiny} \textsuperscript{(63)} ukazují devět hodin a pět minut.

%%K t01-06-2 p
\VUgras{Přestávka} \textsuperscript{(76)} právě skončila a vyučování začíná znovu.
Přestávka se koná na \VUgras{dvoře} \textsuperscript{(75)}. Vidíme několik žáků, kteří si
hrají. \VUgras{Dozorce} \textsuperscript{(74)} na ně dohlíží.

%%K t01-07-1 p
7. --- Na dvoře vidíme ještě různé osoby, totiž \VUgras{inspektora} \textsuperscript{(73)},
\VUgras{ředitele} \textsuperscript{(71)} a \VUgras{cenzora} \textsuperscript{(72)}.

%%K t01-07-2 p
Inspektor kontroluje školy, ředitel řídí lyceum, cenzor dohlíží na studium.

%%K t01-07-3 p
Čtvrtou osobou je \VUgras{vrátný} \textsuperscript{(77)}. Nese pod paží knihu, aby zapsal
nepřítomné.

%%K t01-08-1 p
8. --- Vzadu na dvoře je \VUgras{tělocvičné nářadí} \textsuperscript{(70)} a dílna pro
\VUgras{ruční práce} \textsuperscript{(69)}.

%%K t01-08-2 p
Vpravo od tabulky zahlédneme \VUgras{učebny} \VUgras{hudby} a \VUgras{kreslení}.

%%K t01-noto-1 noto
\VUnotes{143.2mm}{%
(*) U \textit{křestních jmen} se rovněž doporučuje přepisovat je podle latinského tvaru. Je to
jediný způsob, jak uniknout nesčetným obměnám. Ostatně jak vybrat mezi \textit{Giovanni, Jean,
Johann, Jan, Hans, John, Ivan} atd.? Vytvoříme snad zvláštní slovník křestních jmen? Vezměme
si z latiny jejich nominativ a řekněme:
\VUgras{Ioannes, Ioanna; Iulius, Iulia; Iakobus; Andreas; Lukas.} (Úplná
mluvnice).%
}

%%K t01-tit-3 sub
\VUpk{10.2pt}{40}{Podrobný popis.}
\VUsaut{2.4mm}

%%K t01-tit-4 sub
\VUpk{10.2pt}{40}{Dobří žáci.}
\VUsaut{3.0mm}

%%K t01-09-1 p
9. --- Žáci v první lavici vpravo od katedry jsou \VUgras{Jindřich} (*),
\VUgras{Štěpán} a \VUgras{Karel.}

%%K t01-09-2 p
Jindřich je velmi pozorný a pracovitý; poslouchá učitele. Na svém stole má
\VUgras{sešit} \textsuperscript{(28)}, \VUgras{knihu} \textsuperscript{(29)},
\VUgras{slovník} \textsuperscript{(30)} a \VUgras{penál} \textsuperscript{(31)}. Jeho
kniha má mnoho \VUgras{stran}; každá kapitola obsahuje několik \VUgras{odstavců} a každý
\VUgras{řádek} mnoho \VUgras{slov.}

%%K t01-09-4 p
Jeho soused Štěpán \textsuperscript{(24)} je horlivý. Zapisuje si do sešitu úkol na příště. Má
krásné \VUgras{písmo.} Jeho kniha je zavřená. Vidíme jen \VUgras{desky}
\textsuperscript{(27)}. Vedle něho leží jeho \VUgras{kružítko} \textsuperscript{(26)}.

%%K t01-09-5 p
Karel \textsuperscript{(23)} drží knihu v ruce; otevírá ji, aby si ještě jednou přečetl úkol.
Má, jako každý žák, před sebou \VUgras{kalamář} \textsuperscript{(25)}. Je poslušný.

%%K t01-10-1 p
10. --- U vedlejšího stolu sedí \VUgras{Ludvík, Antonín} a \VUgras{Pavel.} Ludvík odříkává
svou \VUgras{lekci} \textsuperscript{(22)} a odpovídá na \VUgras{otázky} učitele. Antonín
\textsuperscript{(21)} je velmi nepozorný; někdy ho učitel dokonce potrestá. Pavel
\textsuperscript{(20)} je dobrý \VUgras{kamarád,} laskavý a ochotný. Je velmi pozorný.

%%K t01-tit-5 sub
\VUsaut{4.2mm}
\VUpk{10.2pt}{40}{Špatní žáci.}
\VUsaut{3.2mm}

%%K t01-11-1 p
11. --- Za Jindřichem sedí \VUgras{Jiří} \textsuperscript{(38)}. Není to špatný chlapec, ale
je \VUgras{upovídaný,} nepozorný a neposedný. Jeho \VUgras{lehkomyslnost} a jeho
\VUgras{neposlušnost} působí \VUgras{zmatek} při vyučování a často si zaslouží přísné
\VUgras{napomenutí.} Jeho \VUgras{sešit na koncepty} \textsuperscript{(35)} je špinavý;
\VUgras{dělá na něj inkoustové skvrny} svým \VUgras{perem} \textsuperscript{(34)}, a proto
stále používá svou \VUgras{gumu} \textsuperscript{(36)}; jeho \VUgras{aktovka}
\textsuperscript{(33)} je roztržená, tak je nedbalý. Proč si nosí své \VUgras{plnicí pero}
\textsuperscript{(37)} do třídy
\VUgras{živých} jazyků?

%%K t01-11-2 p
Jeho soused Edmund \textsuperscript{(39)} ho nemá rád. Je sice poněkud lenivý, ale je
mlčenlivý a klidný. Nežvaní; místo aby poslouchal, čte raději \VUgras{povídky} ve své
\VUgras{čítance}.

%%K t01-11-3 p
Třetím žákem v této lavici je \VUgras{Ludvík} \textsuperscript{(40)}. Je pilný. Píše na bílý
\VUgras{list papíru.} Před sebe si položil \VUgras{savý papír} \textsuperscript{(41)}.

%%K t01-12-1 p
12. --- Žáci ve druhé lavici, vlevo od učitele, jsou velmi mladí. První, malý
\VUgras{Emil,} ještě neumí příliš dobře psát; nosí si svou \VUgras{břidlicovou
tabulku} \textsuperscript{(42)}, své \VUgras{pisátko} a svou \VUgras{houbičku}
\textsuperscript{(43)}.

%%K t01-12-2 p
Vedle něho jsou \VUgras{August} a \VUgras{Bernard} \textsuperscript{(44)}. Tento poslední
pracuje dobře, ale jeho \VUgras{oblečení} je velmi zanedbané.

%%K t01-13-1 p
13. --- Za nimi sedí tři \VUgras{lenoši. Albert} se neučí svým lekcím.
\VUgras{Jakub} \textsuperscript{(47)} spí místo toho, aby poslouchal. Jeho
\VUgras{lenost} mu vynáší \VUgras{výtky} a dokonce \VUgras{tresty.} Konečně
\VUgras{Kristián} \textsuperscript{(48)} je příliš lehkovážný. \VUgras{Chová} se
špatně a nijak se nesnaží polepšit.

%%K t01-14-1 p
14. --- Jeho sousedé se naopak chovají dobře. \VUgras{Arnošt} \textsuperscript{(51)} má rád
pořádek a přesnost; stále používá své
\VUgras{pravítko} \textsuperscript{(50)} a svou \VUgras{tužku}
\textsuperscript{(49)}. Je \VUgras{externí žák.}

%%K t01-14-2 p
\VUgras{František} \textsuperscript{(52)} je \VUgras{internátní žák}; nosí uniformu, jí a
spí v lyceu. Je dobrý \VUgras{kamarád.}

%%K t01-14-3 p
Mořic si ořezává svou \VUgras{tužku} svým \VUgras{nožíkem} \textsuperscript{(56)}; je velmi
pečlivý.

%%K t01-14-4 p
Nosí si všechny své knihy, svou \VUgras{mluvnici} \textsuperscript{(55)}, svůj
\VUgras{zápisník} \textsuperscript{(54)} a dokonce i \VUgras{psací podložku}
\textsuperscript{(53)}. Jeho \VUgras{školní brašna} \textsuperscript{(57)} leží na zemi za
ním.

%%K t01-14-5 p
Oba stoly vzadu ve třídě zaujímají \VUgras{dobří žáci} \textsuperscript{(19)}.

%%K t01-tit-6 sub
\VUsaut{4.6mm}
\VUpk{10.2pt}{40}{Kreslení a hudba.}
\VUsaut{3.4mm}

%%K t01-15-1 p
15. --- V \VUgras{kreslírně} visí na stěně půvabné \VUgras{předlohy} a od stropu je zavěšena
\VUgras{lampa} \textsuperscript{(62)} se svým \VUgras{stínítkem}.
\VUgras{Učitel} \textsuperscript{(60)} je velmi trpělivý; opravuje \VUgras{kresby}
\textsuperscript{(61)} žáků a učí je kreslit \VUgras{bustu} \textsuperscript{(59)} podle
přírody.

%%K t01-16-1 p
16. --- \VUgras{Hudební učebna} se zdá velmi zajímavá. Učí se v ní číst
\VUgras{noty,} zpívat podle not \VUgras{tóny stupnice} \textsuperscript{(67)}
zapsané na \VUgras{osnově} (do, re, mi, fa, sol, la, si\,=\,C. D. E. F. G. A. B.), poznávat
\VUgras{klíče} a zpívat půvabné
\VUgras{melodie.} Noty se kladou na \VUgras{pultík} \textsuperscript{(65)}. Jeden ze
žáků \VUgras{hraje na klavír} \textsuperscript{(64)} a \VUgras{učitel hudby}
\textsuperscript{(66)} \VUgras{udává takt} \textsuperscript{(68)}.

%%K t01-17-1 p
17. --- Zahlédneme kout dílny pro \VUgras{ruční práce} \textsuperscript{(69)}, kde se chlapci
mohou naučit hoblovat dřevo a pilovat železo. To není v každém lyceu.

%%K t01-tit-7 sub
\VUsaut{5.0mm}
\VUpk{10.2pt}{40}{Učebna přírodních věd.}
\VUsaut{3.6mm}

%%K t01-18-1 p
18. --- Učitel matematiky učí \VUgras{žáky} \VUgras{geometrii, počtům, počítání} a
\VUgras{metrické soustavě}. Na tabulce vidíme hlavní geometrické obrazce:
\VUgras{kruh} \textsuperscript{(a)}, \VUgras{čtverec} \textsuperscript{(b)},
\VUgras{trojúhelník} \textsuperscript{(c)}, \VUgras{přímku} \textsuperscript{(d)}.

%%K t01-18-2 p
Vidíme také jeden \VUgras{příklad;} není to \VUgras{sčítání,} ani \VUgras{odčítání,} ani
\VUgras{násobení}: je to opravdu \VUgras{dělení} \textsuperscript{(e)}.

%%K t01-19-1 p
19. --- Na tabulce metrické soustavy vidíme: \VUgras{metr} \textsuperscript{(a)},
\VUgras{váhy} \textsuperscript{(b)} a jejich \VUgras{závaží,} \VUgras{litr}
\textsuperscript{(e)}, \VUgras{dekalitr} \textsuperscript{(f)}, \VUgras{decilitr} a
\VUgras{krychlový metr} \textsuperscript{(d)}.

%%K t01-19-2 p
Žáci studují také \VUgras{přírodopis} \textsuperscript{(11)}, zoologii \textsuperscript{(a)},
botaniku \textsuperscript{(b)}, geologii \textsuperscript{(c)}, a \VUgras{dějepis}
\textsuperscript{(12)}.

%%K t01-19-3 p
Ve skříni jsou přístroje pro \VUgras{fyziku} \textsuperscript{(15)} a pro
\VUgras{chemii} \textsuperscript{(16)}.

%%K t01-19-4 p
Na skříni jsou: \VUgras{koule} \textsuperscript{(14)}, \VUgras{kužel} a
\VUgras{zeměkoule} \textsuperscript{(13)}.

%%K t01-19-5 p
Nad tabulkami visí \VUgras{mapa} znázorňující pět dílů světa: \VUgras{Evropu}
\textsuperscript{(g)}, \VUgras{Asii} \textsuperscript{(e)}, \VUgras{Afriku}
\textsuperscript{(f)}, \VUgras{Ameriku} \textsuperscript{(ab)} a \VUgras{Oceánii}
\textsuperscript{(h)}, a oba velké oceány, \VUgras{Tichý} \textsuperscript{(c)} a
\VUgras{Atlantský} \textsuperscript{(d)}.
\VUsaut{6.0mm}
\VUfilet{22mm}
